% GENERATED FILE. Edit canonical lessons, then run npm run generate:books.
% canonical-source-hash: fnv1a64:e95d070710163b73
% canonical-lessons: TE-C225-akukura, TE-C225-palakura, TE-C225-boppayi, TE-C225-puccakaya, TE-C225-jidipappu, TE-R225-first-pass-words-from-chapters-217-220, TE-R225-second-pass-words-from-chapters-221-225

\chapter{Leafy greens, Spinach, Papaya, Watermelon, Cashew nuts}
\label{ch:te-leafy-greens-spinach-papaya-watermelon-cashew-nuts}
\hlchaptermodality{225}

\begin{chapteropening}
\textbf{By the end of this chapter:} \emph{I can say leafy greens, spinach, a papaya, a watermelon and cashew nuts.}

Complete the last lesson of chapter 225: I can say leafy greens, spinach, a papaya, a watermelon and cashew nuts.
\end{chapteropening}

\section[\texorpdfstring{\te{ఆకుకూర} (ākukūra)}{ఆకుకూర (ākukūra)}]{\te{ఆకుకూర} (ākukūra) — leafy greens}

\label{lesson:TE-C225-akukura}



\begin{quote}
Before the new one: say the Telugu for a mattress, then the Telugu for flour.
\end{quote}

\subsection*{You'll want to know: \te{ఆకుకూర}}
\textbf{\te{ఆకుకూర}} — \emph{ākukūra} — \textquotedblleft{}leafy greens\textquotedblright{}.

\begin{grammarlens}[title={What you've built}]
One more word for this chapter.
\end{grammarlens}

\subsection*{Your turn}
\begin{itemize}
\raggedright
  \item \emph{Say it:} \emph{ākukūra}
  \item \emph{Say it:} \emph{ākukūra}, once more
  \item \emph{Say it:} say \emph{ākukūra}
  \item \emph{Recall:} say the Telugu for a mattress, then the Telugu for flour, then say \emph{ākukūra} again
\end{itemize}

\begin{tcolorbox}[breakable,colback=teal!4,colframe=teal!35!black,title={Before you move on}]
What does \te{ఆకుకూర} mean? (\textquotedblleft{}Leafy greens\textquotedblright{}.) Say it once more.
\end{tcolorbox}
\section[\texorpdfstring{\te{పాలకూర} (pālakūra)}{పాలకూర (pālakūra)}]{\te{పాలకూర} (pālakūra) — spinach}

\label{lesson:TE-C225-palakura}



\begin{quote}
Before the new one: say the Telugu for flour, then the Telugu for leafy greens.
\end{quote}

\subsection*{You'll want to know: \te{పాలకూర}}
\textbf{\te{పాలకూర}} — \emph{pālakūra} — \textquotedblleft{}spinach\textquotedblright{}.

\begin{grammarlens}[title={What you've built}]
One more word for this chapter.
\end{grammarlens}

\subsection*{Your turn}
\begin{itemize}
\raggedright
  \item \emph{Say it:} \emph{pālakūra}
  \item \emph{Say it:} \emph{pālakūra}, once more
  \item \emph{Say it:} say \emph{pālakūra}
  \item \emph{Recall:} say the Telugu for flour, then the Telugu for leafy greens, then say \emph{pālakūra} again
\end{itemize}

\begin{tcolorbox}[breakable,colback=teal!4,colframe=teal!35!black,title={Before you move on}]
What does \te{పాలకూర} mean? (\textquotedblleft{}Spinach\textquotedblright{}.) Say it once more.
\end{tcolorbox}
\section[\texorpdfstring{\te{బొప్పాయి} (boppāyi)}{బొప్పాయి (boppāyi)}]{\te{బొప్పాయి} (boppāyi) — a papaya}

\label{lesson:TE-C225-boppayi}



\begin{quote}
Before the new one: say the Telugu for leafy greens, then the Telugu for spinach.
\end{quote}

\subsection*{You'll want to know: \te{బొప్పాయి}}
\textbf{\te{బొప్పాయి}} — \emph{boppāyi} — \textquotedblleft{}a papaya\textquotedblright{}.

\begin{grammarlens}[title={What you've built}]
One more word for this chapter.
\end{grammarlens}

\subsection*{Your turn}
\begin{itemize}
\raggedright
  \item \emph{Say it:} \emph{boppāyi}
  \item \emph{Say it:} \emph{boppāyi}, once more
  \item \emph{Say it:} say \emph{boppāyi}
  \item \emph{Recall:} say the Telugu for leafy greens, then the Telugu for spinach, then say \emph{boppāyi} again
\end{itemize}

\begin{tcolorbox}[breakable,colback=teal!4,colframe=teal!35!black,title={Before you move on}]
What does \te{బొప్పాయి} mean? (\textquotedblleft{}A papaya\textquotedblright{}.) Say it once more.
\end{tcolorbox}
\section[\texorpdfstring{\te{పుచ్చకాయ} (puccakāya)}{పుచ్చకాయ (puccakāya)}]{\te{పుచ్చకాయ} (puccakāya) — a watermelon}

\label{lesson:TE-C225-puccakaya}



\begin{quote}
Before the new one: say the Telugu for spinach, then the Telugu for a papaya.
\end{quote}

\subsection*{You'll want to know: \te{పుచ్చకాయ}}
\textbf{\te{పుచ్చకాయ}} — \emph{puccakāya} — \textquotedblleft{}a watermelon\textquotedblright{}.

\begin{grammarlens}[title={What you've built}]
One more word for this chapter.
\end{grammarlens}

\subsection*{Your turn}
\begin{itemize}
\raggedright
  \item \emph{Say it:} \emph{puccakāya}
  \item \emph{Say it:} \emph{puccakāya}, once more
  \item \emph{Say it:} say \emph{puccakāya}
  \item \emph{Recall:} say the Telugu for spinach, then the Telugu for a papaya, then say \emph{puccakāya} again
\end{itemize}

\begin{tcolorbox}[breakable,colback=teal!4,colframe=teal!35!black,title={Before you move on}]
What does \te{పుచ్చకాయ} mean? (\textquotedblleft{}A watermelon\textquotedblright{}.) Say it once more.
\end{tcolorbox}
\section[\texorpdfstring{\te{జీడిపప్పు} (jīḍipappu)}{జీడిపప్పు (jīḍipappu)}]{\te{జీడిపప్పు} (jīḍipappu) — cashew nuts}

\label{lesson:TE-C225-jidipappu}



\begin{quote}
Before the new one: say the Telugu for a papaya, then the Telugu for a watermelon.
\end{quote}

\subsection*{You'll want to know: \te{జీడిపప్పు}}
\textbf{\te{జీడిపప్పు}} — \emph{jīḍipappu} — \textquotedblleft{}cashew nuts\textquotedblright{}.

\begin{grammarlens}[title={What you've built}]
One more word for this chapter.
\end{grammarlens}

\subsection*{Your turn}
\begin{itemize}
\raggedright
  \item \emph{Say it:} \emph{jīḍipappu}
  \item \emph{Say it:} \emph{jīḍipappu}, once more
  \item \emph{Say it:} say \emph{jīḍipappu}
  \item \emph{Recall:} say the Telugu for a papaya, then the Telugu for a watermelon, then say \emph{jīḍipappu} again
\end{itemize}

\begin{tcolorbox}[breakable,colback=teal!4,colframe=teal!35!black,title={Before you move on}]
What does \te{జీడిపప్పు} mean? (\textquotedblleft{}Cashew nuts\textquotedblright{}.) Say it once more.
\end{tcolorbox}
\section[(dialogue)]{Words from chapters 217-220 — a first pass}

\label{lesson:TE-R225-first-pass-words-from-chapters-217-220}



\begin{quote}
Say the words for a watermelon and cashew nuts.
\end{quote}

\subsection*{Your turn}
Say these aloud:

\begin{itemize}
\raggedright
  \item together, except, a tradition, a custom, a habit
  \item a result, a goal, an example, a difference, effort
  \item blame, a basis, proof, a development, motivation
  \item free, full, fat, thin, raw
\end{itemize}

\begin{tcolorbox}[breakable,colback=teal!4,colframe=teal!35!black,title={Before you move on}]
Together? (\textbf{\te{కలిసి}}, \emph{kalisi}.) Except? (\textbf{\te{తప్ప}}, \emph{tappa}.) A tradition? (\textbf{\te{సంప్రదాయం}}, \emph{sampradāyaṁ}.) A custom? (\textbf{\te{ఆచారం}}, \emph{ācāraṁ}.) A habit? (\textbf{\te{అలవాటు}}, \emph{alavāṭu}.) A result? (\textbf{\te{ఫలితం}}, \emph{phalitaṁ}.) A goal? (\textbf{\te{లక్ష్యం}}, \emph{lakṣyaṁ}.) An example? (\textbf{\te{ఉదాహరణ}}, \emph{udāharaṇa}.) A difference? (\textbf{\te{తేడా}}, \emph{tēḍā}.) Effort? (\textbf{\te{శ్రమ}}, \emph{śrama}.) Blame? (\textbf{\te{నింద}}, \emph{ninda}.) A basis? (\textbf{\te{ఆధారం}}, \emph{ādhāraṁ}.) Proof? (\textbf{\te{రుజువు}}, \emph{rujuvu}.) A development? (\textbf{\te{పరిణామం}}, \emph{pariṇāmaṁ}.) Motivation? (\textbf{\te{ప్రేరణ}}, \emph{prēraṇa}.) Free? (\textbf{\te{ఉచితం}}, \emph{ucitaṁ}.) Full? (\textbf{\te{నిండు}}, \emph{niṇḍu}.) Fat? (\textbf{\te{లావు}}, \emph{lāvu}.) Thin? (\textbf{\te{సన్నని}}, \emph{sannani}.) Raw? (\textbf{\te{పచ్చి}}, \emph{pacci}.)
\end{tcolorbox}
\section[(dialogue)]{Words from chapters 221-225 — a second pass}

\label{lesson:TE-R225-second-pass-words-from-chapters-221-225}



\begin{quote}
Say the words for brown and silver.
\end{quote}

\subsection*{Your turn}
Say these aloud:

\begin{itemize}
\raggedright
  \item brown, silver, pale, round, sharp
  \item fun, modern, clear, different, fair
  \item personal, private, useless, messy, rude
  \item a ladle, rubbish, the floor, a mattress, flour
  \item leafy greens, spinach, a papaya, a watermelon, cashew nuts
\end{itemize}

\begin{tcolorbox}[breakable,colback=teal!4,colframe=teal!35!black,title={Before you move on}]
Brown? (\textbf{\te{గోధుమ} \te{రంగు}}, \emph{gōdhuma raṅgu}.) Silver? (\textbf{\te{వెండి}}, \emph{veṇḍi}.) Pale? (\textbf{\te{లేత}}, \emph{lēta}.) Round? (\textbf{\te{గుండ్రని}}, \emph{guṇḍrani}.) Sharp? (\textbf{\te{పదునైన}}, \emph{padunaina}.) Fun? (\textbf{\te{సరదా}}, \emph{saradā}.) Modern? (\textbf{\te{ఆధునిక}}, \emph{ādhunika}.) Clear? (\textbf{\te{స్పష్టమైన}}, \emph{spaṣṭamaina}.) Different? (\textbf{\te{భిన్నమైన}}, \emph{bhinnamaina}.) Fair? (\textbf{\te{న్యాయమైన}}, \emph{nyāyamaina}.) Personal? (\textbf{\te{వ్యక్తిగత}}, \emph{vyaktigata}.) Private? (\textbf{\te{ప్రైవేటు}}, \emph{praivēṭu}.) Useless? (\textbf{\te{పనికిరాని}}, \emph{panikirāni}.) Messy? (\textbf{\te{చిందరవందర}}, \emph{cindaravandara}.) Rude? (\textbf{\te{మొరటు}}, \emph{moraṭu}.) A ladle? (\textbf{\te{గరిటె}}, \emph{gariṭe}.) Rubbish? (\textbf{\te{చెత్త}}, \emph{cetta}.) The floor? (\textbf{\te{నేల}}, \emph{nēla}.) A mattress? (\textbf{\te{పరుపు}}, \emph{parupu}.) Flour? (\textbf{\te{పిండి}}, \emph{piṇḍi}.) Leafy greens? (\textbf{\te{ఆకుకూర}}, \emph{ākukūra}.) Spinach? (\textbf{\te{పాలకూర}}, \emph{pālakūra}.) A papaya? (\textbf{\te{బొప్పాయి}}, \emph{boppāyi}.) A watermelon? (\textbf{\te{పుచ్చకాయ}}, \emph{puccakāya}.) Cashew nuts? (\textbf{\te{జీడిపప్పు}}, \emph{jīḍipappu}.)
\end{tcolorbox}
